\documentclass{article}
\usepackage{kotex}
\usepackage[a4paper, top=1.5cm, bottom=1.5cm, left=1.8cm, right=1.8cm, footskip=0.7cm]{geometry}
\usepackage{float}
\usepackage{array}
\usepackage{tablefootnote}

\newcolumntype{C}[1]{>{\centering\arraybackslash}m{#1}}

\title{26-2 `통사론(001)'}
\author{언어학과 이승헌}

\begin{document}
\maketitle
\begin{table}[H]
    \caption{26-2 통사론 스터디 강의계획서}
    \begin{center}
    \renewcommand{\arraystretch}{2.2}
    \begin{tabular}{|C{1.8cm}|C{2.7cm}|C{3.5cm}|C{2cm}|} \hline
    스터디 주차 & 리딩 범위 & 수업 진도\tablefootnote{2025-2 `통사론(001)' 강의계획서를 참조하였음.} & 비고 \\ \hline
    1주차(7/14) & Chapter 1 - 2 & week 1 & \\ \hline
    2주차(7/22) & Chapter 3 - 4 & week 2 - 3 & \\ \hline
    3주차(7/29) & Chapter 5 & week 4 & Part 1 복습 \\ \hline
    4주차\newline(8/4) & Chapter 6 - 7 & week 5, 7\tablefootnote{week 6: 공휴일 휴강} & \\ \hline
    5주차(8/12) & Chapter 8 - 9 & week 8 - 9 & \\ \hline
    6주차(8/19) & Chapter 10 - 11 & week 10 - 11 & Part 2 복습 \\ \hline
    7주차(8/26) & Chapter 12 - 13 & week 12 - 13 & \\ \hline
    8주차\newline(9/2) & Chapter 14 - 15 & week 14 - 15 & Part 3 복습\tablefootnote{리딩한 부분까지만} \\ \hline
    \end{tabular}
    \end{center}
\end{table}
\end{document}